\section*{Bab \ref{ch:b2:planet-formation-meteorites} --- Pembentukan Planet dan Meteorit}

\begin{solution}{exo:b2:planet-formation-meteorites:1}
Gas: \qty{2200}{g.cm^{-2}} pada \qty{1}{au}; pada \qty{5.2}{au},
$2200\times5.2^{-3/2} = 2200/11.86 = \qty{186}{g.cm^{-2}}$. Zat padat:
$33\times0.33 = \qty{10.9}{g.cm^{-2}}$ pada \qty{1}{au};
$33/11.86 = \qty{2.78}{g.cm^{-2}}$ pada \qty{5.2}{au}. Pada \qty{5.2}{au} gas
$186/2.78 = 67$ kali lebih berat daripada zat padat.
\end{solution}

\begin{solution}{exo:b2:planet-formation-meteorites:2}
Aluminium (\qty{1653}{K}, ultrarefraktori), magnesium (\qty{1336}{K}, umum),
besi (\qty{1334}{K}, umum), natrium (\qty{958}{K}, cukup volatil), belerang
(\qty{664}{K}, volatil), oksigen sebagai es air (\qty{180}{K}, sangat volatil).
CAI, kondensat pertama, diperkaya aluminium (serta kalsium dan titanium);
Bumi terkuras unsur-unsur volatilnya, yaitu natrium, belerang, dan air, relatif terhadap
magnesium dan aluminium.
\end{solution}

\begin{solution}{exo:b2:planet-formation-meteorites:3}
\omterm{def:b2:planet-formation-meteorites:fall-find}{Meteorit jatuhan}: \omterm{def:b2:planet-formation-meteorites:chondrule}{kondrit} biasa $856/1107 = \qty{77.3}{\percent}$, \omterm{def:b2:planet-formation-meteorites:classes}{meteorit besi}
$49/1107 = \qty{4.4}{\percent}$. \omterm{def:b2:planet-formation-meteorites:fall-find}{Meteorit temuan}: $38647/44534 = \qty{86.8}{\percent}$ dan
$1020/44534 = \qty{2.3}{\percent}$. \omterm{def:b2:planet-formation-meteorites:fall-find}{Meteorit temuan} Antarktika: $22096/44534 =
\qty{49.6}{\percent}$. Sampel \omterm{def:b2:planet-formation-meteorites:fall-find}{meteorit jatuhan}: setiap kejatuhan yang teramati dikumpulkan apa pun
kelasnya, sedangkan \omterm{def:b2:planet-formation-meteorites:fall-find}{meteorit temuan} bergantung pada seberapa baik \omterm{def:b2:planet-formation-meteorites:fall-find}{meteorit} itu bertahan di tanah,
rupanya, dan tempat orang mencarinya, dan separuh \omterm{def:b2:planet-formation-meteorites:fall-find}{meteorit temuan} berasal dari satu jenis
tempat.
\end{solution}

\begin{solution}{exo:b2:planet-formation-meteorites:4}
$\Delta t = 1.035\ln(\num{5.23e-5}/\num{1.0e-5}) = 1.035\ln 5.23 =
1.035\times1.654 = \qty{1.71}{Ma}$ setelah CAI.
\end{solution}

\begin{solution}{exo:b2:planet-formation-meteorites:5}
$4\pi\Sigma_1 r_1^2 = 4\pi\times\num{2.2e4}\times(\num{1.496e11})^2 =
\qty{6.19e27}{kg}$. Di dalam \qty{30}{au}: $\times\sqrt{30} = 5.48$, sehingga
\qty{3.39e28}{kg} $= 0.017$ massa Matahari; di dalam \qty{100}{au}: $\times10$,
\qty{6.19e28}{kg} $= 0.031$ massa Matahari. Keduanya terletak dalam rentang yang teramati,
lebih dekat ke ujung atasnya: nebula itu adalah cakram biasa.
\end{solution}

\begin{solution}{exo:b2:planet-formation-meteorites:6}
$r = (120/T)^{7/3}$ au: untuk air, $(120/182)^{7/3} = 0.659^{7/3} =
\qty{0.38}{au}$; untuk metana, $(120/41)^{7/3} = 2.93^{7/3} = \qty{12.3}{au}$.
Butiran dalam cahaya Matahari langsung adalah lapisan permukaan cakram; bidang tengahnya,
yang terlindung oleh debu di atasnya, lebih dingin, sehingga es bertahan lebih dekat ke dalam.
\omterm{def:b2:planet-formation-meteorites:planetesimal}{Planetesimal} terbentuk di bidang tengah, sehingga taksiran bidang tengah yang berlaku, asalkan
cakram itu pasif; pemanasan oleh akresi menggeser \omterm{def:b2:planet-formation-meteorites:snow-line}{garis salju} ke luar.
\end{solution}

\begin{solution}{exo:b2:planet-formation-meteorites:7}
Pada \qty{1}{au}: $r^2 = \num{2.238e22}$, $\Sigma_p = \qty{109}{kg.m^{-2}}$, sehingga
$2\pi b r^2\Sigma_p = 31.42\times\num{2.238e22}\times109 = \num{7.66e25}$;
$(\num{7.66e25})^{3/2} = \num{6.71e38}$, dibagi
$(3\times\num{1.989e30})^{1/2} = \num{2.443e15}$: \qty{2.75e23}{kg}
$= 0.046$ massa Bumi. Pada \qty{5.2}{au}: $r^2 = \num{6.05e23}$,
$\Sigma_p = \qty{27.8}{kg.m^{-2}}$, $2\pi b r^2\Sigma_p = \num{5.29e26}$,
$M = \num{1.22e40}/\num{2.443e15} = \qty{4.98e24}{kg} = 0.83$ massa Bumi.
Dengan zat padat \num{3.5} kali lebih banyak, $M$ tumbuh sebanding dengan $3.5^{3/2} = 6.55$: 5.5 massa
Bumi.
\end{solution}

\begin{solution}{exo:b2:planet-formation-meteorites:8}
Per kilogram terdapat $\num{1.01e19}$ atom, sehingga
$\num{1.01e19}\times\num{5.00e-13} = \qty{5.05e6}{J.kg^{-1}}$. Kenaikan
suhunya akan menjadi $\num{5.05e6}/837 = \qty{6030}{K}$. Menghangatkan dari \num{250} hingga
\qty{1700}{K} memerlukan $837\times1450 = \qty{1.21e6}{J.kg^{-1}}$: aluminium-26
memasok sekitar empat kali lipatnya.
\end{solution}

\begin{solution}{exo:b2:planet-formation-meteorites:9}
$\num{5.07e10}/\num{7.49e7} = 677$. Umur rata-rata besi-60 adalah
$2.62/\ln2 = \qty{3.78}{Ma}$. Daya awal per gram:
$\num{5.07e10}/1.035 = \num{4.90e10}$ dan $\num{7.49e7}/3.78 = \num{1.98e7}$ erg per
gram per juta tahun, perbandingan sekitar \num{2470}: aluminium-26 jauh lebih dominan.
\end{solution}

\begin{solution}{exo:b2:planet-formation-meteorites:10}
Kedua daya itu sama ketika $\num{4.90e10}\,e^{-t/1.035} = \num{1.98e7}\,e^{-t/3.78}$,
yaitu $t(1/1.035 - 1/3.78) = \ln(\num{4.90e10}/\num{1.98e7}) = \ln 2475 = 7.81$;
dengan $0.966 - 0.265 = 0.702$ per juta tahun, $t = \qty{11.1}{Ma}$. Pada saat itu
besi-60 tinggal menyisakan $\num{7.49e7}\,e^{-11.1/3.78} = \num{7.49e7}\times0.053 =
\qty{4.0e6}{erg.g^{-1}}$, yaitu \qty{400}{J.kg^{-1}}: ribuan kali terlalu
sedikit untuk melebur apa pun.
\end{solution}

\begin{solution}{exo:b2:planet-formation-meteorites:11}
$\tau = 1.035\times\num{3.156e13} = \qty{3.267e13}{s}$ dan
$H_0\tau = \qty{6.21e6}{J.kg^{-1}}$. Untuk \qty{1200}{K}:
$t_a^* = 1.035\ln\bigl(\num{6.21e6}/(837\times950)\bigr) = 1.035\ln7.81 =
\qty{2.13}{Ma}$; untuk \qty{1700}{K}: $1.035\ln\bigl(\num{6.21e6}/(837\times1450)\bigr)
= 1.035\ln5.12 = \qty{1.69}{Ma}$. Membagi dua $H_0$ mengurangi kedua waktu itu dengan $\tau\ln2 =
\qty{0.72}{Ma}$, karena $t_a^* = \tau\ln H_0 + \tau\ln(\tau/c(T_c -
T_0))$ dan hanya suku pertama yang berubah: pergeserannya satu waktu paruh, berapa pun
suhu kritisnya.
\end{solution}

\begin{solution}{exo:b2:planet-formation-meteorites:12}
$\kappa = 2.1/(3300\times837) = \qty{7.60e-7}{m^2.s^{-1}}$ dan
$\sqrt{\kappa\tau} = \sqrt{\num{7.60e-7}\times\num{3.267e13}} = \qty{4980}{m}$,
sekitar \qty{5}{km}. Pusat benda berdiameter \qty{100}{km} berada \qty{50}{km} di bawah permukaan,
sepuluh kali jarak ini: panasnya tidak dapat lolos selama aluminium-26 meluruh, dan
proposisi itu berlaku. Benda berdiameter \qty{2}{km} kehilangan panasnya secepat ia
menghasilkannya. Peleburan memerlukan jari-jari beberapa kali \qty{5}{km}, yaitu diameter sekitar
\qtyrange{10}{20}{km} atau lebih, seperti yang didapati model-model.
\end{solution}

\begin{solution}{pb:b2:planet-formation-meteorites:1}
\textbf{1.} \omterm{def:b2:planet-formation-meteorites:fall-find}{Meteorit jatuhan}: kedatangannya terlihat dan batu itu dipungut keesokan paginya.
Batu itu belum terlapukkan, dan menambah sampel yang tidak berat sebelah dari apa yang mencapai Bumi.

\textbf{2.} $856/1107 = 0.77$: sekitar tiga peluang dari empat.

\textbf{3.} \omterm{def:b2:planet-formation-meteorites:chondrule}{Kondrul} berdiameter \qty{0.3}{mm}, logam \qty{9}{\percent}, dan olivin dengan
\qty{18}{\percent} fayalit: \omterm{def:b2:planet-formation-meteorites:chondrule}{kondrit} H (logam \qty{10}{\percent}, Fa 16--20).

\textbf{4.} \omterm{def:b2:planet-formation-meteorites:chondrule}{Kondrul} yang kabur dan mineral yang seragam menandakan metamorfisme termal: tipe
5 hingga 6. \omterm{def:b2:planet-formation-meteorites:chondrule}{Kondrit} biasa mencapai paling tinggi sekitar \qty{1220}{K}, yaitu
\qty{947}{\degreeCelsius}.

\textbf{5.} \omterm{def:b2:planet-formation-meteorites:chondrule}{Kondrit} biasa hanya mengandung \qtyrange{0.1}{1}{\percent} CAI menurut
volume. Batu itu tersusun atas \omterm{def:b2:planet-formation-meteorites:chondrule}{kondrul} dan matriks yang terhimpun setelah CAI, yang
merupakan zat padat tertua: batuan apa pun yang mengandungnya lebih muda.

\textbf{6.} Setiap mineral menerima saat pembentukannya magnesium yang sama ditambah aluminium-26
sebanding dengan aluminiumnya; aluminium-26 itu sejak itu telah menjadi
magnesium-26, sehingga makin banyak aluminium per magnesium yang dimiliki sebuah mineral, makin banyak
magnesium-26 yang diperolehnya.

\textbf{7.} $(0.139590 - 0.139500)/(15.0 - 0.05) = \num{0.000090}/14.95 =
\num{6.02e-6}$: perbandingan $\ce{^{26}Al}/\ce{^{27}Al}$ \omterm{def:b2:planet-formation-meteorites:chondrule}{kondrul} itu ketika
terbentuk.

\textbf{8.} $0.139500 + \num{6.02e-6}\times3.00 = 0.139518$: terletak pada garis itu.

\textbf{9.} $\Delta t = 1.035\ln(\num{5.23e-5}/\num{6.02e-6}) = 1.035\ln8.69 =
1.035\times2.162 = \qty{2.24}{Ma}$.

\textbf{10.} $4567.5 - 2.2 = \qty{4565.3}{\Ma}$.

\textbf{11.} $\num{5.23e-5}\,e^{-10/1.035} = \num{5.23e-5}\times\num{6.4e-5} =
\num{3.3e-9}$, magnesium-26 yang jauh terlalu sedikit untuk diukur. Hafnium-182 (waktu paruh
\qty{8.9}{Ma}) atau mangan-53 (\qty{3.7}{Ma}) cocok untuk peristiwa sekitar sepuluh juta tahun
setelah CAI.

\textbf{12.} $\kappa = 2.1/(3300\times837) = \qty{7.60e-7}{m^2.s^{-1}}$;
$R^2/\kappa = \num{4e8}/\num{7.60e-7} = \qty{5.26e14}{s} = \qty{16.7}{Ma}$, enam belas
kali umur rata-rata aluminium-26: pusatnya mempertahankan panasnya.

\textbf{13.} $\num{1.9e-7}\,e^{-2.4/1.035} = \num{1.9e-7}\times0.0984 =
\qty{1.87e-8}{W.kg^{-1}}$.

\textbf{14.} $\num{1.87e-8}\times\num{3.267e13} = \qty{6.11e5}{J.kg^{-1}}$.

\textbf{15.} $\Delta T = \num{6.11e5}/837 = \qty{730}{K}$; suhu puncak
$250 + 730 = \qty{980}{K}$.

\textbf{16.} Di atas \qty{750}{K}, sehingga batuan itu tersinter dan termetamorfosis; di bawah
\qty{1200}{K}, sehingga baik logam maupun sulfida tidak melebur. Batu itu adalah \omterm{def:b2:planet-formation-meteorites:chondrule}{kondrit} yang terpanaskan,
bertipe petrologi tinggi, di bawah \qty{1220}{K} pada pertanyaan 4: konsisten.

\textbf{17.} $1.035\ln\bigl(\num{6.21e6}/(837\times950)\bigr) = \qty{2.13}{Ma}$:
sekitar \qty{0.3}{Ma} lebih awal.

\textbf{18.} Kenaikannya $\Delta T(1 - e^{-(t - t_a)/\tau})$; \qty{90}{\percent}
memerlukan $e^{-(t - t_a)/\tau} = 0.1$, $t - t_a = 1.035\ln10 = \qty{2.38}{Ma}$, yaitu
\qty{4.8}{Ma} setelah CAI.

\textbf{19.} $\Delta T = 7420\,e^{-0.8/1.035} = 7420\times0.462 = \qty{3430}{K}$,
suhu puncak \qty{3680}{K}, jauh di atas \qty{1700}{K} saat silikat cair.
Benda itu melebur (panas peleburan menahan suhunya) dan
terdiferensiasi menjadi inti logam serta mantel dan kerak silikat.

\textbf{20.} $e^{(2.4 - 0.8)/1.035} = e^{1.546} = 4.7$ kali lebih banyak.

\textbf{21.} Seng dan timbal mengembun di dekat \qty{727}{K}, cukup volatil; magnesium
(\qty{1336}{K}) dan aluminium (\qty{1653}{K}) mengembun jauh lebih awal. Zat padat
di daerah batu itu terkumpul sebelum gas cukup dingin agar seluruh seng dan
timbal mengembun.

\textbf{22.} $33\times0.33\times2.5^{-3/2} = 10.89/3.95 = \qty{2.76}{g.cm^{-2}}$,
\qty{27.6}{kg.m^{-2}}.

\textbf{23.} $r^2 = (3.74\times10^{11})^2 = \num{1.40e23}$;
$2\pi\times5\times\num{1.40e23}\times27.6 = \num{1.21e26}$;
$(\num{1.21e26})^{3/2}/\num{2.443e15} = \qty{5.45e23}{kg}$, $0.091$ massa Bumi.
Benda induk itu, $\tfrac43\pi(\num{2e4})^3\times3300 = \qty{1.1e17}{kg}$, lima
juta kali lebih ringan: sebuah \omterm{def:b2:planet-formation-meteorites:planetesimal}{planetesimal}, bukan embrio.

\textbf{24.} Pusat benda induk itu mencapai puncak sekitar
\textbf{\qty[text-series-to-math=true]{707}{\degreeCelsius}} ($980 - 273$): cukup panas untuk
mengkristalkan ulang \omterm{def:b2:planet-formation-meteorites:chondrule}{kondrit} itu, terlalu dingin untuk melebur logamnya.
\end{solution}
